% GENERATED_BY: oc_core_monograph_translation_draft_pipeline_v1.py
% AI_ASSISTED_TRANSLATION_DRAFT: TRUE
% THEOREM_REVIEW_STATUS: PENDING
% THEOREM_REVIEW_MODE: FORMAL_AI_GATE__CODEX_CLI_CHATGPT
% THEOREM_REVIEW_REVIEWER_ID:
% THEOREM_REVIEW_AUTHORITY_CLASS:
% THEOREM_REVIEWED_AT:
% THEOREM_REVIEW_DOSSIER_REF:
% SOURCE_LANGUAGE: EN
% TARGET_LANGUAGE: RU
% SOURCE_EN_REF: content/m_spaces/m3.tex
% ================================================================
% ==== FILE: content/m_spaces/m3.tex
% ================================================================

% ==============================
%  Ontology of Continua — Core
%  Meta-Space: M3
%  FULL MODULE — FINAL VERSION
% ==============================

\subsubsection{\texorpdfstring{$M_3$}{M_3} Обзор}
\label{sec:m3-overview}

$M_3$ — это мета-пространство, задающее условия допустимости для
трёхмерных континуумов $K_3$.
Это первое мета-пространство, в котором:

\begin{itemize}
    \item становится допустимой третья независимая ось $A_3$,
    \item появляется реактивная структура (химическое связывание/диссоциация),
    \item валентностные состояния входят в допустимое конфигурационное пространство,
    \item в трёх измерениях возможны диффузия и транспорт,
    \item локальные графы взаимодействий (реакционные сети) становятся устойчивыми объектами,
    \item пороги активации и потенциалы реакций становятся определёнными,
    \item пространственная неоднородность и градиенты получают динамическую поддержку.
\end{itemize}

$M_3$ является математической и физической средой, в которой
химический континуум $K_3$ может существовать, эволюционировать и проходить
структурные переходы к более высоким континуумам, таким как $K_4$.

% ---------------------------------------------------------------
\subsubsection*{1. Допустимое конфигурационное пространство \texorpdfstring{$\Omega(M_3)$}{\Omega(M_3)}}

Допустимое множество конфигураций в $M_3$:

\[
\Omega(M_3)
=
\Big\{
(\phi,\, G_{\mathrm{chem}},\, \rho,\, g_{ij}) \ \Big|
\phi \in C^{0}(X_3, V),\
G_{\mathrm{chem}} \in \mathcal{G}_{\mathrm{RAF}},\
\rho \in L^1(X_3),\
g_{ij}\in C^{0}
\Big\}.
\]

Ключевые условия допустимости:

\begin{enumerate}
    \item \textbf{Трёхмерная топология.}
    $X_3$ должен быть 3-мерным многообразием, поддерживающим локальные
    координатные карты, дифференцируемость и диффузию.

    \item \textbf{Локальные реакционные сети.}
    $G_{\mathrm{chem}}$ должен быть представим как конечный, локально
    поддерживаемый реакционный граф, удовлетворяющий RAF-условиям замыкания
    и каталитического действия.

    \item \textbf{Поля плотности.}
    $\rho(x)$ задаёт молекулярные или атомные концентрации и должно
    быть интегрируемым с конечной общей массой.

    \item \textbf{Совместимость метрики.}
    $g_{ij}$ описывает пространственную геометрию, но пока не включает
    релятивистскую структуру; на данном уровне допустимы только
    римановы метрики (положительно определённые).

    \item \textbf{Локальность.}
    Все взаимодействия должны быть выражаемы через локальные потенциалы
    и локальные реакционные ядра.
\end{enumerate}

Таким образом, $\Omega(M_3)$ расширяет геометрическую и физическую
допустимость $M_2$ в химический режим.

% ---------------------------------------------------------------
\subsubsection*{2. Допустимые оси \texorpdfstring{$A(M_3)$}{A(M_3)}}

В $M_3$ становится допустимой новая ось $A_3$, для которой выполняется:

\[
A_3 \not\subset \operatorname{span}(A_1, A_2),
\qquad
T_2 \ge \Theta_{\mathrm{dim}}.
\]

Дополнительные допустимые внутренние оси:

\begin{itemize}
    \item $A_{\mathrm{val}}$: состояния валентности (возможности связи),
    \item $A_{\mathrm{react}}$: допустимые режимы реакций,
    \item $A_{\mathrm{conf}}$: конфигурационные состояния молекул.
\end{itemize}

Эти внутренние оси кодируют химическую структуру внутри мета-пространства.

% ---------------------------------------------------------------
\subsubsection*{3. Потенциалы \texorpdfstring{$P(M_3)$}{P(M_3)}}

Допустимые потенциалы в $M_3$ соответствуют химическим и пространственным
взаимодействиям:

\[
P(M_3)
=
\left\{
U_{\mathrm{bond}} + U_{\mathrm{rep}} + U_{\mathrm{conf}}
+ U_{\mathrm{diff}} + U_{\mathrm{local}}
\right\}.
\]

где:

\begin{itemize}
    \item $U_{\mathrm{bond}}$ — потенциалы связи с минимумами на
    устойчивых длинах связи,
    \item $U_{\mathrm{rep}}$ — короткодействующие репульсивные потенциалы,
    предотвращающие коллапс,
    \item $U_{\mathrm{conf}}$ — внутренние конфигурационные потенциалы
    для молекулярных состояний,
    \item $U_{\mathrm{diff}}$ — потенциалы, связанные с градиентами концентрации,
    \item $U_{\mathrm{local}}$ — любой локально определённый вклад в энергию.
\end{itemize}

Ограничения:

\begin{enumerate}
    \item потенциалы должны быть ограничены снизу,
    \item потенциалы должны быть почти всюду дифференцируемы,
    \item они должны задавать локальную, а не глобальную или дальнодействующую
    структуру.
\end{enumerate}

% ---------------------------------------------------------------
\subsubsection*{4. Пороги \texorpdfstring{$\Theta(M_3)$}{\Theta(M_3)}}

$M_3$ формализует несколько химических порогов:

\begin{itemize}
    \item $\Theta_{\mathrm{bond}}$ — минимальная энергия для разрыва связи,
    \item $\Theta_{\mathrm{act}}$ — порог активации для реакций,
    \item $\Theta_{\mathrm{cluster}}$ — минимальная энергия устойчивого кластера,
    \item $\Theta_{\mathrm{grad}}$ — порог поддержания пространственных градиентов,
    \item $\Theta_{\mathrm{RAF}}$ — порог замыкания RAF-сети.
\end{itemize}

Стабильность $K_3$ требует:

\[
T_3 < \min\{\Theta_{\mathrm{bond}}, \Theta_{\mathrm{cluster}},
\Theta_{\mathrm{RAF}}\}.
\]

где $T_3$ — структурная натяжка конфигураций в $\Omega(K_3)$.

% ---------------------------------------------------------------
\subsubsection*{5. Потоки \texorpdfstring{$J(M_3)$}{J(M_3)}}

Допустимые потоки включают:

\[
J(M_3) =
\big(
J_{\mathrm{diff}},\
J_{\mathrm{react}},\
J_{\mathrm{grad}},\
\nabla_i T_3
\big).
\]

Интерпретация:

\begin{itemize}
    \item $J_{\mathrm{diff}}$ — диффузионные потоки в 3D,
    \item $J_{\mathrm{react}}$ — локальные реакционные потоки по RAF-правилам,
    \item $J_{\mathrm{grad}}$ — потоки поддержания или ослабления
    локальных градиентов,
    \item $\nabla_i T_3$ — потоки, обусловленные градиентами структурной натяжки.
\end{itemize}

Все потоки должны сохранять локальный закон сохранения массы.

% ---------------------------------------------------------------
\subsubsection*{6. Циклы и реакционные сети}

$M_3$ — это первое мета-пространство, где допустимы химические циклы:

\[
C_{\mathrm{chem}}
=
\big\{
\text{связывание} \to
\text{преобразование} \to
\text{разрыв-связи}
\to \dots
\big\}.
\]

RAF-сети становятся допустимыми циклами, когда:

\[
E_{\mathrm{in}} \ge \Theta_{\mathrm{act}}
\quad\text{и}\quad
\text{условия замыкания выполнены}.
\]

Эти циклы образуют минимальные строительные блоки возникновения
биологических континуумов ($K_4$).

% ---------------------------------------------------------------
\subsubsection*{7. Граница \texorpdfstring{$\partial\Omega(M_3)$}{\partial\Omega(M_3)} и коллапс}

Конфигурация приближается к границе $\partial\Omega(M_3)$, если:

\begin{itemize}
    \item потенциалы связи становятся сингулярными,
    \item реакционный поток расходится,
    \item градиенты становятся неинтегрируемыми ($|\nabla\rho| \to \infty$),
    \item теряется RAF-замыкание ($G_{\mathrm{chem}} \notin \mathcal{G}_{\mathrm{RAF}}$),
    \item суммарная структурная натяжка превышает порог:
    \[
    T_3 \ge \Theta_3.
    \]
\end{itemize}

Пересечение $\partial\Omega(M_3)$ разрушает химическую непрерывность и
приводит к коллапсу $K_3$.

% ---------------------------------------------------------------
\subsubsection*{8. Связь с \texorpdfstring{$M_2$}{M_2} и \texorpdfstring{$M_4$}{M_4}}

Переход $K_2 \to K_3$ становится допустимым, когда:

\[
A_3 \in A(M_3),
\qquad
\Omega(K_3) \subseteq \Omega(M_3),
\qquad
T_2 \ge \Theta_{\mathrm{dim}}.
\]

Переход к $K_4$ требует:

\[
\Omega(K_4) \subseteq \Omega(M_4),
\qquad
\partial\Omega(M_3) \ \text{обеспечивает ограниченную проницаемость},
\qquad
\Theta_{\mathrm{mem}} > 0.
\]

Таким образом, $M_3$ — мета-пространство:

\begin{itemize}
    \item химической структуры,
    \item реакционных сетей,
    \item трёхмерной диффузии,
    \item валентности и связывания,
    \item порогов активации,
    \item устойчивых градиентов.
\end{itemize}

Это структурная основа, обеспечивающая существование химических континуумов $K_3$
и возникновение протоклеточных структур в $K_4$.
